\section{Вычислительный юмор как исследовательская проблема}

Юмор давно рассматривается как сложный случай для обработки естественного языка, поскольку успешные шутки зависят не только от грамматической правильности. Они также зависят от несоответствия ожиданиям, общего знания, тайминга и ожиданий аудитории. Ранние обзоры уже описывали юмор как одну из наиболее сложных целей для искусственного интеллекта, отмечая, что реализованные системы в основном ограничивались узкими механизмами каламбуров, а более широкая обработка юмора потребовала бы значительного знания о мире и способности к рассуждению \citep{ritchie2001current}. Два более поздних обзора приходят к сходному выводу с разных сторон. \citet{amin2020survey} показывают, что генерация вычислительного юмора исторически обменивала широту на контролируемость: шаблонные и основанный на правилах системы могли хорошо работать при наличии явного юмористического механизма, но только в сильно ограниченных доменах. Более недавно \citet{Loakman2025} и \citet{lemmens2026computational} утверждают, что эпоха больших языковых моделей расширила покрытие, но не решила центральную проблему, поскольку беглый текст не эквивалентен юмору, который читатели действительно находят смешным.

Дополнительная трудность связана с теоретическим плюрализмом. Теории юмора предлагают несколько повторяющихся классов объяснений---прежде всего подходы, основанные на несоответствие, но также традиции превосходство и разрядка,---и при этом ни одна рамка не исчерпывает феномен полностью. Это важно вычислительно, потому что разные типы шуток опираются на разные механизмы. Например, \citet{kao2016computational} формализуют несоответствие в каламбурах как неоднозначность плюс отличительность и показывают, что эти переменные предсказывают человеческие оценки смешного. Их работа важна не только тем, что находит применение классической теории юмора, но и тем, что показывает пределы узкой формализации: модель, объясняющая омофонические каламбуры, не начинает автоматически объяснять нарративные шутки, тематический шутки или социально ситуированный юмор. Современная основанный на теории работа по обнаружению юмора исходит из той же предпосылки: единой всеобъемлющей теории юмора недостаточно, а вычислительные модели обычно усваивают только фрагменты юмора \citep{demarez2024thinc}. Для данной работы следствие является методологическим: юмор следует рассматривать как творческую задачу со множеством потенциальных механизмов, а не как один скалярный атрибут коммуникации.

\section{Эволюция подходов к генерации юмора}

До появления больших языковых моделей системы генерации юмора обычно были узкими по замыслу. Как резюмирует \citet{amin2020survey}, ранние работы часто опирались на вручную написанные шаблоны, лексические ресурсы, фонетические преобразования или жестко ограниченные процедуры генерации. У этих систем было важное преимущество: они кодировали по крайней мере один явный юмористический механизм. Их недостатком была низкая адаптивность. Они генерировали юмор только при тщательно выбранных лингвистических условиях и поэтому плохо масштабировались за пределы каламбуров, устойчивых форм шуток или ограниченных доменов.

Большие языковые модели переворачивают этот компромисс. Они могут генерировать открытый юмористический текст по гораздо более широкому диапазону запросы, но беглость не гарантирует устойчивую юмористическую компетентность. Это видно в недавних крупномасштабных оценках. \citet{Zhang2024HumorAI} вводят мультимодальный контрольный набор для подписей к карикатурам, построенный более чем на 250 миллионах человеческих оценок для более чем 2.2 миллиона подписи, и показывают, что современные методы дообучение остаются ограниченными на творческих задачах. Даже сильные системы, такие как GPT-4 и Claude, остаются ниже лучших человеческих участников. Исследования понимания юмора указывают в том же направлении. \citet{xu2024good} показывают, что большие языковые модели демонстрируют то, что они называют lazy каламбур генерация: вместо слияния двух значений в один компактный каламбур модель часто отдельно называет оба связанных слова или значения, поверхностно удовлетворяя задаче. Этот результат методологически важен, поскольку показывает, что простые лексический успех метрики могут переоценивать способность к каламбурам. Оценка должна также отличать структурное качество и оригинальность от механически собранной или скопированной игры слов. \citet{cocchieri2025you} строят устойчивый к загрязнение данных контрольный набор юмора и показывают, что большинство языковых моделей все еще уступают людям в понимании, разрешении и генерации каламбуров. Выходя за пределы каламбуров, \citet{loakman2025comparing} показывают, что современные модели, включая модели рассуждения, не объясняют надежно все типы шуток, что подчеркивает, насколько большая часть литературы по-прежнему сфокусирована на относительно простых формах юмора.

В совокупности эти исследования показывают, что современные модели лучше воспроизводят поверхностную форму знакомых шуток, чем обобщают юмористические механизмы. Это различие важно для данной работы. Система может хорошо генерировать текст, похожий на шутку, но не быть хорошей в юморе в более сильном смысле, необходимом для открытой генерации.

\section{Проблемы оценки юмора}

Оценка является центральной трудностью вычислительного юмора. Юмор субъективен, но проблема не только в субъективности в разговорном смысле. То, что считается хорошей шуткой, зависит от аудитории, контекста, подачи и оцениваемого измерения. Даже статьи о набора данныхх подчеркивают, что юмор меняется по теме, времени и получателю, из-за чего трудно выделить одну целевую переменную \citep{weller2020rjokes}. Поэтому современные работы по юмору все чаще избегают трактовки смешность как скалярной метрики.

Наиболее ясный пример дает \citet{Sakabe2025}, оценивающий Oogiri ответы по шести измерениям: новизна, ясность, релевантность, интеллект, эмпатия и общая смешность. Их результаты показывают, что современные языковые модели генерируют ответы между низким и средним уровнем человеческих результатов, но заметно отстают по эмпатия. Авторы также утверждают, что этот дефицит помогает объяснить, почему модельный оценка плохо воспроизводит человеческую оценку юмора. Значение этого результата двоякое. Во-первых, он показывает, что качество юмора не сводится к вопросу ``является this смешной?'' Во-вторых, он показывает, что человеческие и модельные суждения могут отдавать приоритет разным качествам.

Дополнительная перспектива представлена у \citet{winters2025evaluating}, которые сравнивают человеческих комиков и GPT-4 в живом импровизационном постановкае, а не в офлайн только текстовый контрольный набор. Их исследование ценно тем, что задает более сильную точку сравнения: и люди, и модель должны быстро отвечать на одни и те же предложения аудитории. Человеческие шутки все еще предпочитаются немного чаще, но разрыв не является подавляющим, а дизайн подчеркивает, насколько сильно подача и ситуативный контекст влияют на оценку юмора. Это полезная коррекция к мышлению только через контрольный наборs. Текст, хорошо работающий в статической оценке, не обязательно является текстом, который лучше всего работает в реальном комедийном взаимодействии.

Эти результаты также уточняют роль методос помощью языковой модели-оценщика. Структурированные оценочные рамки, такие как у \citet{Hashemi2024}, полезны, потому что заменяют одномерный оценивание явными рубриками и калибровкой. Однако эта работа также показывает, что сырые LLM рубрика суждения не автоматически хорошо согласуются с человеческими рейтингами. Для генерации юмора это означает, что LLM оценщикs лучше рассматривать как измерительные инструменты с ограниченной областью применимости, а не как прозрачную замену человеческих предпочтений.

\section{Перспектива рассуждения для генерации юмора}

Один из наиболее важных сдвигов в современных исследованиях генерации юмора заключается в переходе от с одним примером продолжение к явному рассуждению или поэтапному творческому поиску. На уровне построение подсказок \citet{chen2023prompt} моделируют генерацию юмора через пошаговые инструкции комедийного письма для GPT-3. \citet{zhong2024let} идут дальше и утверждают, что стандартный цепочки рассуждений построение подсказок плохо подходит творческим задачам; они предлагают парадигму Leap--Thought для Oogiri-style генерации юмора, где акцент делается на ассоциативных скачках, а не на строго последовательной декомпозиции.

Более явные системы генерации усиливают это направление. \citet{Tikhonov2024} реконструируют написание короткая шутка шутки как многошаговый процесс и сообщают об улучшениях относительно без примеров GPT-4 и других базовые модели в человеческой оценке. \citet{Kim2025} описывают генерацию юмора как комбинацию когнитивных, социальных и творческих навыков. В их исследовании meme-подписи к изображениям дополненный навыками системы существенно превосходят базовые модели на больших языковых моделях и приближаются к лучшим человеческим подписи для целевой аудитории. В мультимодальных условиях \citet{ZhangJ2025} предлагают направляемый теорией multi-этап рассуждение рамка, направленный на улучшение и интерпретируемости, и качества юмора.

Эти статьи важны, потому что коллективно отвергают предположение, что юмор лучше всего порождается мгновенным завершением. Вместо этого они моделируют генерацию юмора как процесс поиска по предпосылки, ассоциации, переосмысления и зависящий от аудитории варианты, аналогично тому, как человеческие эксперты подходят к комедийному письму или другим творческим задачам. Между предпосылка и шуткой может существовать информационный разрыв, который заполняется только во время инференс и может быть латентным. В то же время эти подходы разделяют ограничение, напрямую релевантное данной работе: структура рассуждения задается жестко. Число этапов, тип промежуточных представлений и порядок операций фиксируются исследователем. Поэтому такие системы являются сильной демонстрацией полезности рассуждения, но более слабым свидетельством того, что модель может сама научиться, когда и как рассуждать для юмора.

\section{Стандартное постобучение плохо подходит юмору}

Самый прямой путь к улучшению генерации юмора --- с учителем дообучение на больших корпусах шуток. На практике, однако, такая целевая функция плохо соответствует структуре задачи. Набор данныхы шуток содержат множество повторяющихся завязки, знакомых паттернов игра слов и конвенциональных форм панчлайн. В результате с учителем дообучение может улучшать стилистическую правдоподобность, не обязательно улучшая творческое обобщение. Крупные ресурсы юмора, такие как rJokes, ценны масштабом и тематическим разнообразием, но они не были построены как с несколькими эталонными ответами генерация наборы данных, где много разных ответы одновременно могут быть правильными \citep{weller2020rjokes}.

Современная литература усиливает это опасение с двух сторон. Во-первых, юмор-specific работы показывают, что даже очень большие предпочтение наборы данных не делают творческий выравнивание простой задачей. \citet{Zhang2024HumorAI} прямо подчеркивают ограничения RLHF- и DPO-style дообучение на творческом подписи к изображениям. Во-вторых, работы по после-обучение вне юмора делают различие запоминание--обобщение более резким. \citet{Chu2025} показывают, что обучение с подкреплением обобщает эффективнее, чем с учителем дообучение, на ненаблюдаемый текстовый и визуальный варианты, тогда как с учителем дообучение склонен запоминать обучение паттерны и испытывать трудности out распределение. Этот контраст особенно важен для генерации юмора, где слабый запрос может иметь много хороших продолжений, а запоминание частых оболочек шуток не равно созданию уместного и неожиданного панчлайн.

Preference оптимизация имеет собственные трудности. RLHF предполагает, что награда модель может выучить устойчивый приближенная мера для целевого поведения, но суждения о юморе шумны и неоднородны. \citet{Gao2023} в более общем виде показывают, что агрессивная оптимизация несовершенной награда модель может ухудшать истинный качество. Humor-specific оценка показывает, что риск усиливается в творческих доменах: \citet{Sakabe2025} показывают систематическое расхождение между человеческий и модель юмор суждения, а \citet{winters2025evaluating} показывают, что качество юмора зависит от аспектов качество, которые трудно захватить офлайн текстовый награда сигналы. По этим причинам ни с учителем дообучение, ни стандартная предпочтение оптимизация не решают центральную проблему чисто.

\section{Постобучение, ориентированное на рассуждение}

Недавний успех ориентированный на рассуждение после-обучение в основном пришел из доменов с объективной верификацией. \citet{Wei2022} показывают, что промежуточное рассуждение может существенно менять поведение модели на задачах, требующих многошагового инференс. \citet{Lightman2023} расширяют эту логику от построение подсказок к разметка, показывая, что процесс разметка превосходит только по результату разметка на контрольный набор MATH. \citet{DeepSeek2025} затем демонстрируют, что обучение с подкреплением может вызывать длинный формат рассуждение без опоры на написанные людьми цепочки thought, причем наиболее ясные выигрыши наблюдаются в верифицируемых доменах, таких как математика, программирование и связанные STEM задачи.

Это создает центральную перенос проблема для юмора. Юмор также может выигрывать от расширенного рассуждения, но обычно не имеет детерминированного верификатора, способного объявить шутку правильной. Современные без внешнего верификатора подходы для доменов вне математики и программирования пытаются закрыть этот разрыв, заменяя явный проверяемую награду внутренними вероятностями, назначенными эталонные ответы. \citet{Zhou2025} предлагают VeriFree и показывают, что при с одним эталонным ответом предположение максимизация условной вероятности эталонный ответ при заданных запрос и траектория рассуждения эквивалентна максимизации точное совпадение верификатор целевая функция в ожидании. \citet{Yu2025} развивают близкий подход RLPR, где награда формулируется как среднее правдоподобие по отдельным токенам для повышения стабильности дисперсии.

Однако для юмора с одним эталонным ответом предположение не является малой технической деталью; это неверная модель данных. Prompts вроде ``write a joke about a frog'' по природе допускают много правдоподобных ответов, и любой набор данных будет наблюдать только небольшое и смещенное их подмножество. Иными словами, юмор одновременно является с несколькими эталонными ответами и только частично наблюдаемый задачей. Это означает, что без внешнего верификатора рассуждение методы концептуально релевантны генерации юмора, но не применимы напрямую без переформулировки награда вокруг множеств эталонных ответов, а не одиночных целей.

\section{Исследовательский пробел}

Рассмотренная литература оставляет точный пробел. Исследования генерации юмора все чаще показывают, что поэтапный рассуждение, навык декомпозиция или направляемый теорией генерация могут улучшать ответы, но такие системы в основном опираются на жесткие инференс процессы \citep{chen2023prompt,zhong2024let,Tikhonov2024,Kim2025,ZhangJ2025}. Работы по ориентированный на рассуждение обучение с подкреплением показывают, что полезные рассуждение поведение можно вызывать обучением, но сильнейшие существующие формулировки предполагают либо объективную верифицируемость, либо один эталонный ответ \citep{Lightman2023,DeepSeek2025,Zhou2025,Yu2025}. В то же время исследования оценки юмора показывают, что смешность является многомерной, контекстно-зависимой и лишь несовершенно захватывается модель оценщикs \citep{Zhang2024HumorAI,Sakabe2025,winters2025evaluating,Hashemi2024}.

Данная работа находится на пересечении этих результатов. Она утверждает, что юмор выигрывает от рассуждения, но отвергает необходимость задавать фиксированный творческий сценарий. Она также развивает без внешнего верификатора обучение с подкреплением, но утверждает, что творческие задачи требуют с несколькими эталонными ответами формулировка, поскольку несколько ответы могут быть одновременно хорошими, а наблюдаемый набор эталонных ответов неполон. Поэтому исследовательский вопрос состоит в том, можно ли адаптировать без внешнего верификатора рассуждение целевая функцияs к частично наблюдаемый с несколькими эталонными ответами юмор генерация так, чтобы поощрять творческий поиск без схлопывания в запомненный обучение выборки.
